\chapter{Especificación detallada de requisitos}
\label{cap:especificacion-detallada}

Este capítulo desarrolla los flujos de los casos de uso y los requisitos
de información y calidad recogidos en la especificación de origen. Se
conservan los códigos para mantener la trazabilidad entre documentos.

\textbf{Nota de alcance.} Las referencias a re-votación o sustitución (RF-11,
estados “Sustituido” y “re-votación permitida”) proceden de la especificación
v1.0 y no están implementadas en el MVP (D-03). El estado “Sustituido” se
conserva en el modelo como referencia de diseño para RF-11 (P9), pero no es
operativo: la urna registra un único voto aceptado por token de emisión y no
genera ese estado.

\section{Casos de uso}
\label{sec:casos-uso-detallados}

\subsection{RF-01: Gestionar organizaciones}

\textbf{Objetivos asociados:} OBJ-1\par
\textbf{Requisitos asociados:} RI-1\par
\textbf{Descripción:} El sistema deberá comportarse como se describe a continuación cuando un Superadministrador (ACT-1) da de alta una organización.\par
\textbf{Precondición:} El usuario está autenticado (RF-04) como Superadministrador.\par
\textbf{Flujo normal:}
\begin{enumerate}
\item[1] El Superadministrador solicita dar de alta una organización
\item[2] El Sistema solicita el nombre, identificador y administrador inicial de la organización.
\item[3] El Superadministrador introduce los datos solicitados.
\item[4] El Sistema comprueba que el identificador no está en uso.
\item[5] El Sistema crea la organización con su espacio de datos aislado y asigna el rol de Administrador de organización al administrador inicial.
\end{enumerate}
\textbf{Postcondición:} La organización existe con su espacio de datos aislado y el administrador inicial tiene asignado el rol de Administrador de organización.\par
\textbf{Excepciones y alternativas:}
\begin{itemize}
\item 4: Si el identificador ya está en uso, el Sistema informa de la situación y solicita introducir otro identificador.
\end{itemize}
\textbf{Comentarios:} Ninguno\par

\subsection{RF-02: Gestionar roles y permisos}

\textbf{Objetivos asociados:} OBJ-1\par
\textbf{Requisitos asociados:} RI-2\par
\textbf{Descripción:} El sistema deberá comportarse como se describe a continuación cuando un Administrador de organización (ACT-2) asigna o cambia el rol de un usuario de su organización (el Superadministrador (ACT-1) puede hacerlo en cualquier organización).\par
\textbf{Precondición:} El usuario está autenticado y el usuario destino pertenece a la organización.\par
\textbf{Flujo normal:}
\begin{enumerate}
\item[1] El Administrador solicita gestionar usuarios y roles
\item[2] El Sistema muestra los usuarios de la organización con su rol
\item[3] El Administrador selecciona un usuario o solicita crear uno nuevo
\item[4] El Sistema muestra los roles asignables
\item[5] El Administrador selecciona el rol
\item[6] El Sistema guarda el rol y aplica sus permisos desde la siguiente acción del usuario
\end{enumerate}
\textbf{Postcondición:} El usuario tiene el nuevo rol y sus permisos efectivos.\par
\textbf{Excepciones y alternativas:}
\begin{itemize}
\item 5: Si se intenta dejar la organización sin Administrador, el Sistema lo impide e informa
\item 5: Si el actor no puede asignar ese rol, el Sistema deniega la operación
\end{itemize}
\textbf{Comentarios:} Prioridad M. Se asume que solo un Superadministrador puede asignar el rol de Superadministrador.\par

\subsection{RF-03: Gestionar censo}

\textbf{Objetivos asociados:} OBJ-2\par
\textbf{Requisitos asociados:} RI-3, RNF-4\par
\textbf{Descripción:} El sistema deberá comportarse como se describe a continuación cuando un Administrador da de alta, importa o modifica votantes del censo.\par
\textbf{Precondición:} El Gestor de censo está autenticado.\par
\textbf{Flujo normal:}
\begin{enumerate}
\item[1] El Gestor solicita gestionar el censo
\item[2] El Sistema muestra los censos existentes
\item[3] El Gestor elige la acción
\item[]\textit{Si elige alta manual:}
\item[3.1] El Sistema solicita nombre, email o teléfono
\item[3.2] El Gestor introduce los datos
\item[]\textit{Si elige importar CSV:}
\item[3.1] El Sistema solicita el fichero CSV
\item[3.2] El Gestor selecciona el fichero
\item[3.3] El Sistema valida el fichero y muestra un resumen (altas, duplicados, errores)
\item[3.4] El Gestor confirma la importación
\item[4] El Sistema detecta duplicados y no los da de alta
\item[5] El Sistema guarda los cambios e informa del resultado
\end{enumerate}
\textbf{Postcondición:} El censo está actualizado y sin duplicados.\par
\textbf{Excepciones y alternativas:}
\begin{itemize}
\item 3.3: Si el fichero no es válido, el Sistema informa de las líneas erróneas y permite descartarlas o cancelar
\item 4: Si hay duplicados, el Sistema los omite e informa de cuáles
\item 5: Si el censo pertenece a un proceso abierto, el Sistema rechaza altas y bajas e informa
\end{itemize}
\textbf{Comentarios:} Prioridad M. Se asume que el censo queda congelado mientras el proceso está abierto.\par

\subsection{RF-04: Autenticar usuario de gestión}

\textbf{Objetivos asociados:} OBJ-1\par
\textbf{Requisitos asociados:} RI-2\par
\textbf{Descripción:} El sistema deberá comportarse como se describe a continuación cuando un Superadministrador, Administrador, Gestor de censo o Auditor desea iniciar sesión.\par
\textbf{Precondición:} La cuenta existe y está activa.\par
\textbf{Flujo normal:}
\begin{enumerate}
\item[1] El Usuario solicita iniciar sesión
\item[2] El Sistema solicita usuario y contraseña
\item[3] El Usuario introduce sus credenciales
\item[4] El Sistema valida las credenciales
\item[5] El Sistema solicita el segundo factor
\item[6] El Usuario introduce el segundo factor
\item[7] El Sistema valida el segundo factor
\item[8] El Sistema inicia la sesión y muestra la pantalla de inicio de su rol
\end{enumerate}
\textbf{Postcondición:} Sesión iniciada con el rol y la organización correspondientes.\par
\textbf{Excepciones y alternativas:}
\begin{itemize}
\item 4: Si las credenciales son incorrectas, el Sistema informa del error y permite reintentar.
\item 7: Si el código del segundo factor es incorrecto o ha caducado, el Sistema informa del error y permite reintentar o solicitar un nuevo código.
\end{itemize}
\textbf{Comentarios:} Prioridad M. La autenticación se realiza mediante el mecanismo configurado para la plataforma. Como opción de diseño, podrá integrarse un proveedor de identidad compatible con OIDC/SSO (RF-05); no se exige una cuenta institucional. El segundo factor podrá realizarse mediante un código enviado al correo electrónico del usuario.\par

\subsection{RF-05: Iniciar sesión (RF-05a gestión y RF-05b votante)}

\textbf{Objetivos asociados:} OBJ-1, OBJ-3\par
\textbf{Requisitos asociados:} RI-2, RI-3\par
\textbf{Descripción:} El sistema deberá permitir que los perfiles de gestión (RF-05a) y los votantes (RF-05b) inicien sesión mediante el mecanismo de autenticación habilitado para la plataforma o el proceso. La integración con un proveedor de identidad (ACT-6), por ejemplo mediante OIDC/SSO, es una opción de diseño.\par
\textbf{Precondición:} La cuenta está activa o, en el caso del votante, pertenece al censo del proceso. Si se usa un proveedor externo, este está configurado y disponible.\par
\textbf{Flujo normal:}
\begin{enumerate}
\item[1] El usuario solicita iniciar sesión y elige uno de los mecanismos habilitados
\item[2] El Sistema solicita las credenciales o inicia el flujo del proveedor de identidad, si está configurado
\item[3] El usuario se autentica con el mecanismo seleccionado
\item[4] El Sistema valida la respuesta y comprueba la cuenta, el rol y, cuando corresponda, la pertenencia al censo
\item[5] El Sistema inicia la sesión o concede el acceso al proceso
\end{enumerate}
\textbf{Postcondición:} Sesión iniciada con el rol y el ámbito de organización o proceso correspondientes.\par
\textbf{Excepciones y alternativas:}
\begin{itemize}
\item 3: Si el usuario cancela o no se autentica, el Sistema vuelve a la pantalla de inicio
\item 2--4: Si el proveedor no responde, las credenciales son incorrectas o la respuesta no es válida, el Sistema informa y permite reintentar o usar otro mecanismo habilitado
\item 4: Si la cuenta no tiene un rol autorizado o el votante no figura en el censo, el Sistema deniega el acceso e informa
\end{itemize}
\textbf{Comentarios:} Prioridad Should. El requisito describe el inicio de sesión de forma independiente del proveedor y del protocolo. OIDC/SSO es una alternativa de diseño que deberá concretarse si se adopta; no implica SSO institucional obligatorio.\par

\subsection{RF-06: Crear y configurar proceso de votación}

\textbf{Objetivos asociados:} OBJ-2\par
\textbf{Requisitos asociados:} RI-4, RI-5, RNF-4\par
\textbf{Descripción:} El sistema deberá comportarse como se describe a continuación cuando un Administrador de organización (ACT-2) crea y configura un proceso de votación, incluyendo la definición del tipo de voto y de la papeleta.\par
\textbf{Precondición:} El Administrador está autenticado y existe un censo (RF-03).\par
\textbf{Flujo normal:}
\begin{enumerate}
\item[1] El Administrador solicita crear un proceso
\item[2] El Sistema solicita título y descripción
\item[3] El Administrador introduce el título y la descripción.
\item[4] El Sistema solicita el calendario (apertura y cierre)
\item[5] El Administrador introduce fechas y horas
\item[6] El Sistema solicita el censo participante.
\item[7] El Administrador selecciona el censo participante.
\item[8] El Sistema solicita definir la papeleta.
\item[9] El Sistema muestra los tipos de voto disponibles: opción única y sí/no/abstención.
\item[10] El Administrador selecciona el tipo de voto.
\item[11] El Administrador introduce las opciones y decide si habilita el voto en blanco.
\end{enumerate}
\textbf{Postcondición:} El proceso de votación existe en estado Borrador con su papeleta definida.\par
\textbf{Excepciones y alternativas:}
\begin{itemize}
\item 12: Si existen datos incompletos o incoherentes en el proceso o en la papeleta, el Sistema informa del error y permite corregir los datos sin guardar.
\end{itemize}
\textbf{Comentarios:} Prioridad M. Incluye la definición del tipo de voto y de la papeleta.\par

\subsection{RF-07: Gestionar ciclo de vida del proceso}

\textbf{Objetivos asociados:} OBJ-2\par
\textbf{Requisitos asociados:} RI-4, RI-9\par
\textbf{Descripción:} El sistema deberá comportarse como se describe a continuación cuando un Administrador de organización (ACT-2) programa y archiva un proceso, y el Sistema lo abre y lo cierra automáticamente en las fechas fijadas.\par
\textbf{Precondición:} El proceso existe y el Administrador está autenticado.\par
\textbf{Flujo normal:}
\begin{enumerate}
\item[1] El Administrador solicita programar el proceso
\item[2] El Sistema valida que el proceso está completo.
\item[3] El Sistema cambia el proceso a Programado
\item[4] Al llegar la fecha de apertura, el Sistema cambia el proceso a Abierto y habilita el acceso de los votantes (RF-08)
\item[5] Al llegar la fecha de cierre, el Sistema cambia el proceso a Cerrado y deja de aceptar votos
\item[6] El Sistema publica la huella final de la urna (RF-12), inicia el escrutinio (RF-13) y, al terminar, cambia el proceso a Escrutado
\item[7] El Administrador solicita archivar el proceso
\item[8] El Sistema cambia el proceso a Archivado
\end{enumerate}
\textbf{Postcondición:} Estado del proceso actualizado y registrado (RF-19).\par
\textbf{Excepciones y alternativas:}
\begin{itemize}
\item 2: Si el proceso está incompleto, el Sistema informa y no lo programa.
\end{itemize}
\textbf{Comentarios:} Prioridad M. Lo amplían RF-13 (escrutinio al cerrar) y RF-17 (avisos). No se permite abrir ni cerrar un proceso manualmente de forma anticipada; la apertura y el cierre se realizan automáticamente en las fechas configuradas.\par

\subsection{RF-08: Acceder como votante}

\textbf{Objetivos asociados:} OBJ-3\par
\textbf{Requisitos asociados:} RI-3, RI-12, RNF-1, RNF-2, RNF-3, RNF-5, RNF-9\par
\textbf{Descripción:} El sistema deberá comportarse como se describe a continuación cuando un Votante (ACT-4) accede a la papeleta de un proceso abierto; el código lo envía el Servicio de notificaciones (ACT-7).\par
\textbf{Precondición:} El proceso está Abierto, el votante está en el censo y ha recibido su enlace personal (RF-17).\par
\textbf{Flujo normal:}
\begin{enumerate}
\item[1] El Votante abre su enlace personal
\item[2] El Sistema valida el enlace y que el proceso está abierto
\item[3] El Sistema solicita la verificación adicional del proceso
\item[]\textit{Si es un código OTP:}
\item[3.1] El Sistema envía el código por email
\item[3.2] El Votante introduce el código
\item[3.3] El Sistema valida el código
\item[4] El Sistema concede acceso a la papeleta
\end{enumerate}
\textbf{Postcondición:} El votante está identificado para ese proceso.\par
\textbf{Excepciones y alternativas:}
\begin{itemize}
\item 2: Si el enlace no es válido o ha caducado, el Sistema informa y no concede acceso
\item 2: Si el proceso no está abierto, el Sistema informa de su estado y calendario
\item 3.3: Si el código es incorrecto o ha caducado, el Sistema permite reintentar o reenviarlo.
\end{itemize}
\textbf{Comentarios:} Prioridad M. La verificación adicional configurada (por ejemplo, OTP o autenticación federada) evita que un enlace reenviado permita votar a otra persona. La federación de identidad es una opción de diseño, no un requisito institucional.\par

\subsection{RF-09: Emitir voto}

\textbf{Objetivos asociados:} OBJ-3\par
\textbf{Requisitos asociados:} RI-5, RI-6, RI-7, RNF-1, RNF-2, RNF-3, RNF-5, RNF-7, RNF-9\par
\textbf{Descripción:} El sistema deberá comportarse como se describe a continuación cuando un Votante (ACT-4) emite su voto en un proceso abierto.\par
\textbf{Precondición:} El votante ha accedido (RF-08) y no tiene un voto contabilizado, o el proceso permite re-votar (RF-11).\par
\textbf{Flujo normal:}
\begin{enumerate}
\item[1] El Sistema muestra la papeleta.
\item[2] El Votante selecciona sus opciones
\item[3] El Sistema valida la selección según el tipo de papeleta
\item[4] El Sistema muestra un resumen y pide confirmación
\item[5] El Votante confirma
\item[6] El Sistema comprueba que no hay ya un voto contabilizado del votante
\item[7] El Sistema registra el voto en la urna.
\item[8] El Sistema marca al votante como «ha votado»
\item[9] El Sistema confirma la emisión y ofrece el recibo (RF-10)
\end{enumerate}
\textbf{Postcondición:} El voto está en la urna y el votante figura como votante.\par
\textbf{Excepciones y alternativas:}
\begin{itemize}
\item 3: Si la selección no es válida, el Sistema informa y no avanza
\item 6: Si el votante ya votó y no se permite re-votar, el Sistema informa y no registra nada
\item 7: Si el proceso se cierra mientras vota o falla el registro, el Sistema informa y no queda ningún voto parcial ni duplicado
\end{itemize}
\textbf{Comentarios:} Prioridad M. El voto se registra en la urna y el votante queda marcado como «ha votado». Como máximo se contabiliza un voto por votante.\par

\subsection{RF-10: Obtener y verificar recibo del voto}

\textbf{Objetivos asociados:} OBJ-3, OBJ-4\par
\textbf{Requisitos asociados:} RI-6, RI-8, RNF-1, RNF-2, RNF-3, RNF-6\par
\textbf{Descripción:} El sistema deberá comportarse como se describe a continuación cuando un Votante (ACT-4) quiere comprobar que su voto consta en la urna.\par
\textbf{Precondición:} El votante ha emitido un voto (RF-09).\par
\textbf{Flujo normal:}
\begin{enumerate}
\item[1] El Sistema genera un recibo del voto.
\item[2] El Sistema entrega el recibo al votante
\item[3] El Votante solicita verificar su recibo
\item[4] El Sistema busca el recibo en la urna publicada
\item[5] El Sistema informa de que el voto consta en la urna
\end{enumerate}
\textbf{Postcondición:} Ninguna.\par
\textbf{Excepciones y alternativas:}
\begin{itemize}
\item 4: Si el recibo no se encuentra, el Sistema informa de que el voto no consta
\end{itemize}
\textbf{Comentarios:} Prioridad M. Amplía RF-09. Si el votante vuelve a votar (RF-11), el recibo anterior se marca como sustituido.\par

\subsection{RF-11: Re-votar}

\textbf{Estado en el MVP:} No implementado. D-03 excluye RF-11 del MVP; la exclusión fue confirmada por el tutor el 30/10/2026 mediante correo electrónico, y el caso de uso permanece como trabajo futuro (P9). Se conserva aquí por trazabilidad con la especificación v1.0.\par

\textbf{Objetivos asociados:} OBJ-3\par
\textbf{Requisitos asociados:} RI-6, RI-7, RI-8, RNF-1, RNF-2, RNF-7, RNF-9\par
\textbf{Descripción:} El sistema deberá comportarse como se describe a continuación cuando un Votante (ACT-4) cambia su voto antes del cierre.\par
\textbf{Precondición:} El proceso está Abierto, permite re-votación y el votante ya ha votado y ha accedido (RF-08).\par
\textbf{Flujo normal:}
\begin{enumerate}
\item[1] El Votante solicita cambiar su voto
\item[2] El Sistema informa de que el nuevo voto sustituirá al anterior
\item[3] El Votante confirma
\item[4] El Votante emite el nuevo voto (RF-09)
\item[5] El Sistema registra el nuevo voto y marca el anterior como sustituido sin borrarlo
\item[6] El Sistema entrega el nuevo recibo
\end{enumerate}
\textbf{Postcondición:} Solo cuenta el último voto del votante.\par
\textbf{Excepciones y alternativas:}
\begin{itemize}
\item 3: Si la votante cancela, su voto anterior sigue siendo válido
\item 5: Si el proceso se cierra antes de registrar el nuevo voto, el anterior sigue siendo válido
\end{itemize}
\textbf{Comentarios:} Prioridad W (fuera de esta versión; era C en la especificación v1.0). Amplía RF-15. La re-votación mitiga la coacción: el votante puede cambiar su voto hasta el cierre.\par

\subsection{RF-12: Publicar y verificar registro de auditoría}

\textbf{Objetivos asociados:} OBJ-4\par
\textbf{Requisitos asociados:} RI-9\par
\textbf{Descripción:} El sistema deberá comportarse como se describe a continuación cuando un Auditor (ACT-3) consulta y verifica el registro de la urna, o un Votante (ACT-4) consulta la huella publicada.\par
\textbf{Precondición:} Existe un proceso con eventos en la urna.\par
\textbf{Flujo normal:}
\begin{enumerate}
\item[1] El Auditor solicita ver el registro de auditoría de un proceso
\item[2] El Sistema muestra los eventos de la urna (apertura, votos recibidos sin identidad, sustituciones, cierre) con su huella
\item[3] El Auditor solicita verificar la integridad
\item[4] El Sistema comprueba la integridad del registro.
\item[5] El Sistema informa del resultado de la verificación
\item[]\textit{Si el proceso está cerrado:}
\item[6] El Sistema muestra la huella final publicada, que cualquier Votante puede consultar
\end{enumerate}
\textbf{Postcondición:} Ninguna.\par
\textbf{Excepciones y alternativas:}
\begin{itemize}
\item 2: Si el proceso no tiene eventos, el Sistema informa
\item 5: Si la cadena está rota, el Sistema alerta al Auditor y al Administrador y marca una incidencia de integridad
\end{itemize}
\textbf{Comentarios:} Prioridad S.\par

\subsection{RF-13: Escrutar y verificar resultado}

\textbf{Objetivos asociados:} OBJ-4\par
\textbf{Requisitos asociados:} RI-4, RI-6, RI-10, RNF-6\par
\textbf{Descripción:} El sistema deberá comportarse como se describe a continuación cuando un proceso se cierra y el Sistema lanza automáticamente el escrutinio; después, un Auditor (ACT-3) o un Votante (ACT-4) verifica el resultado.\par
\textbf{Precondición:} El proceso está Cerrado.\par
\textbf{Flujo normal:}
\begin{enumerate}
\item[1] El Sistema lanza el escrutinio al cerrarse el proceso
\item[2] El Sistema obtiene los votos contabilizables (el último de cada votante).
\item[3] El Sistema cuenta los votos por opción.
\item[4] El Sistema calcula la participación.
\item[5] El Sistema guarda el resultado y publica su huella.
\item[6] El Auditor o el Votante solicita verificar el resultado
\item[7] El Sistema muestra los datos públicos necesarios para recalcularlo
\item[8] El Sistema informa de si el resultado coincide
\end{enumerate}
\textbf{Postcondición:} Proceso en estado Escrutado con resultado y huella publicados.\par
\textbf{Excepciones y alternativas:}
\begin{itemize}
\item 5: Si falla el recuento, el Sistema no publica nada, avisa al Administrador y permite reintentar.
\item 8: Si no coincide, el Sistema alerta de una incidencia
\end{itemize}
\textbf{Comentarios:} Prioridad M. Amplía RF-07. Queda por decidir qué significa «público»: cualquier persona o solo los participantes del proceso.\par

\subsection{RF-14: Consultar panel de participación}

\textbf{Objetivos asociados:} OBJ-4\par
\textbf{Requisitos asociados:} RI-3, RI-7\par
\textbf{Descripción:} El sistema deberá comportarse como se describe a continuación cuando un Administrador de organización (ACT-2) consulta la participación de un proceso.\par
\textbf{Precondición:} El Administrador está autenticado.\par
\textbf{Flujo normal:}
\begin{enumerate}
\item[1] El Administrador solicita el panel de participación
\item[2] El Sistema muestra el número de votantes del censo, los votos emitidos y el porcentaje de participación.
\item[3] El Sistema actualiza el panel mientras el proceso está abierto.
\end{enumerate}
\textbf{Postcondición:} Ninguna.\par
\textbf{Excepciones y alternativas:}
\begin{itemize}
\item 2: Si el proceso aún no está abierto, el Sistema muestra solo el censo
\item 2: Si el proceso está cerrado, muestra la participación final
\end{itemize}
\textbf{Comentarios:} Prioridad M. Ninguna vista muestra el recuento por opción antes del cierre.\par

\subsection{RF-15: Consultar y exportar resultados}

\textbf{Objetivos asociados:} OBJ-4\par
\textbf{Requisitos asociados:} RI-10\par
\textbf{Descripción:} El sistema deberá comportarse como se describe a continuación cuando un Administrador de organización (ACT-2), o un Votante (ACT-4) si el proceso lo permite, consulta y exporta los resultados.\par
\textbf{Precondición:} El proceso está Escrutado.\par
\textbf{Flujo normal:}
\begin{enumerate}
\item[1] El usuario solicita ver los resultados
\item[2] El Sistema comprueba que el usuario puede verlos
\item[3] El Sistema muestra los resultados con gráficos y participación.
\item[4] El usuario solicita exportarlos
\item[5] El Sistema genera el fichero en formato CSV y lo entrega.
\end{enumerate}
\textbf{Postcondición:} Ninguna.\par
\textbf{Excepciones y alternativas:}
\begin{itemize}
\item 2: Si el proceso no está escrutado, el Sistema informa de que aún no hay resultados
\item 2: Si el usuario no tiene permiso, el Sistema deniega el acceso
\end{itemize}
\textbf{Comentarios:} Prioridad M. Los resultados pueden ser consultados por el Administrador de organización y por los votantes cuando el proceso lo permita. La exportación de resultados se realiza únicamente en formato CSV.\par

\subsection{RF-16: Generar acta de resultados en PDF}

\textbf{Objetivos asociados:} OBJ-4\par
\textbf{Requisitos asociados:} RI-10\par
\textbf{Descripción:} El sistema deberá comportarse como se describe a continuación cuando un Administrador de organización (ACT-2) genera el acta de un proceso.\par
\textbf{Precondición:} El proceso está Escrutado y el Administrador está autenticado.\par
\textbf{Flujo normal:}
\begin{enumerate}
\item[1] El Administrador solicita el acta
\item[2] El Sistema compone el acta (proceso, calendario, censo, participación, resultado y huellas)
\item[3] El Sistema genera el PDF y lo ofrece para descarga
\item[4] El Administrador lo descarga
\end{enumerate}
\textbf{Postcondición:} Acta disponible en PDF con huella.\par
\textbf{Excepciones y alternativas:}
\begin{itemize}
\item 1: Si el proceso no está escrutado, el Sistema informa de que el acta no está disponible
\item 3: Si falla la generación, el Sistema informa y permite reintentar
\end{itemize}
\textbf{Comentarios:} Prioridad S. Amplía RF-15. El contenido del acta debe validarse con las exigencias de cada organización.\par

\subsection{RF-17: Enviar notificaciones}

\textbf{Objetivos asociados:} OBJ-3\par
\textbf{Requisitos asociados:} RI-11, RI-12\par
\textbf{Descripción:} El sistema deberá comportarse como se describe a continuación cuando se programa, abre o cierra un proceso y el Sistema avisa a los votantes mediante el Servicio de notificaciones (ACT-7). El Administrador de organización (ACT-2) configura las plantillas y los canales.\par
\textbf{Precondición:} El proceso está programado y el Administrador ha configurado plantillas y canales.\par
\textbf{Flujo normal:}
\begin{enumerate}
\item[1] El Sistema inicia el envío de invitaciones al programarse el proceso
\item[2] El Sistema genera el enlace personal de cada votante (RF-08) y compone el mensaje en su idioma
\item[3] El Sistema entrega los mensajes al Servicio de notificaciones
\item[4] El Servicio de notificaciones informa del resultado de entrega
\item[5] El Sistema registra el estado de cada envío
\item[6] En las fechas configuradas, el Sistema envía recordatorios a quien aún no ha votado
\item[7] Al cierre, el Sistema envía el aviso de cierre
\end{enumerate}
\textbf{Postcondición:} Votantes avisados y estado de envío registrado.\par
\textbf{Excepciones y alternativas:}
\begin{itemize}
\item 4: Si una entrega falla, el Sistema reintenta y la marca; el Administrador puede consultar los fallos
\item 2: Si un votante no tiene dato de contacto para el canal, el Sistema lo omite y lo lista.
\end{itemize}
\textbf{Comentarios:} Prioridad M. Amplía RF-08.\par

\subsection{RF-18: Gestionar privacidad de datos (RGPD)}

\textbf{Objetivos asociados:} OBJ-1\par
\textbf{Requisitos asociados:} RI-2, RI-3\par
\textbf{Descripción:} El sistema deberá comportarse como se describe a continuación cuando un Administrador de organización (ACT-2) exporta o borra los datos personales de un interesado a petición de este.\par
\textbf{Precondición:} El Administrador está autenticado.\par
\textbf{Flujo normal:}
\begin{enumerate}
\item[1] El Administrador solicita la gestión de privacidad
\item[]\textit{Si exporta datos de un interesado:}
\item[2.1] El Administrador identifica al interesado.
\item[2.2] El Sistema compila sus datos personales y los exporta en formato legible por máquina.
\item[]\textit{Si borra datos de un interesado:}
\item[3.1] El Administrador identifica al interesado y confirma la solicitud de borrado.
\item[3.2] El Sistema borra o anonimiza sus datos personales.
\end{enumerate}
\textbf{Postcondición:} Los datos del interesado han sido exportados o borrados, según la solicitud.\par
\textbf{Excepciones y alternativas:}
\begin{itemize}
\item 3.2: Si el interesado ya votó, el Sistema anonimiza sus datos personales y conserva el voto y el registro de auditoría.
\end{itemize}
\textbf{Comentarios:} Prioridad S. El voto no es exportable por interesado porque no está vinculado a su identidad. Si el votante ya ha votado, se anonimizan sus datos personales y se conserva el voto.\par

\subsection{RF-19: Consultar trazabilidad de acciones}

\textbf{Objetivos asociados:} OBJ-4\par
\textbf{Requisitos asociados:} RI-13\par
\textbf{Descripción:} El sistema deberá comportarse como se describe a continuación cuando un Auditor (ACT-3) o un Administrador de organización (ACT-2) consulta las acciones administrativas registradas.\par
\textbf{Precondición:} El usuario está autenticado. El Sistema registra automáticamente cada acción administrativa.\par
\textbf{Flujo normal:}
\begin{enumerate}
\item[1] El usuario solicita la trazabilidad
\item[2] El Sistema solicita filtros (fechas, usuario, tipo de acción, proceso)
\item[3] El usuario los introduce
\item[4] El Sistema muestra las acciones: quién, qué, sobre qué objeto y cuándo
\end{enumerate}
\textbf{Postcondición:} Ninguna.\par
\textbf{Excepciones y alternativas:}
\begin{itemize}
\item 4: Si no hay resultados, el Sistema informa
\end{itemize}
\textbf{Comentarios:} Prioridad S. El registro lo hace el Sistema en cada caso de uso de administración y no se puede editar. El Administrador solo ve las acciones de su organización.\par

\section{Requisitos de información}
\label{sec:datos-detallados}

\subsection{RI-1: Información sobre organizaciones}

\textbf{Objetivos asociados:} OBJ-1\par
\textbf{Requisitos asociados:} RF-01\par
\textbf{Descripción:} El sistema deberá almacenar la información correspondiente a las organizaciones (tenants) que utilizan la plataforma.\par
\textbf{Datos específicos:} Identificador de organización: 20 caracteres alfanuméricos (único) Nombre visible: 100 caracteres Estado: Activa o Suspendida Fecha de alta: dd/mm/aaaa hh:mm\par
\textbf{Comentarios:} Los datos de cada organización están aislados lógicamente de los de las demás (arquitectura multi-tenant).\par

\subsection{RI-2: Información sobre usuarios y roles}

\textbf{Objetivos asociados:} OBJ-1\par
\textbf{Requisitos asociados:} RF-01, RF-02, RF-04, RF-05\par
\textbf{Descripción:} El sistema deberá almacenar la información correspondiente a los usuarios de gestión de la plataforma y a su rol.\par
\textbf{Datos específicos:} Identificador de usuario: entero positivo Organización: identificador de organización (RI-1) Nombre y apellidos: 60 caracteres Correo electrónico: 100 caracteres (único por organización) Teléfono: 15 dígitos (opcional) Rol: Superadministrador, Administrador o Auditor Contraseña: se almacena solo su huella (hash) Segundo factor activado: Sí o No Estado de la cuenta: Activa, Bloqueada o Inactiva Fecha del último acceso: dd/mm/aaaa hh:mm\par
\textbf{Comentarios:} Los votantes no son usuarios de gestión: se almacenan en el censo (RI-3). Se asume que un usuario pertenece a una única organización, salvo el Superadministrador, que pertenece a la plataforma.\par

\subsection{RI-3: Información sobre censo}

\textbf{Objetivos asociados:} OBJ-2\par
\textbf{Requisitos asociados:} RF-03, RF-05, RF-08, RF-14, RF-18\par
\textbf{Descripción:} El sistema deberá almacenar la información correspondiente a los votantes de cada organización.\par
\textbf{Datos específicos:} Identificador de votante: entero positivo Organización: identificador de organización (RI-1) Nombre y apellidos: 60 caracteres Correo electrónico: 100 caracteres Teléfono: 15 dígitos (opcional) Identificador externo: 50 caracteres (opcional, solo si el mecanismo de autenticación seleccionado lo requiere)\par
\textbf{Comentarios:} El control de duplicados se hace mediante los datos definidos por la organización, como el correo electrónico; un identificador externo solo se conserva si resulta necesario para el mecanismo configurado.\par

\subsection{RI-4: Información sobre procesos de votación}

\textbf{Objetivos asociados:} OBJ-2\par
\textbf{Requisitos asociados:} RF-06, RF-07, RF-08, RF-09, RF-12, RF-13, RF-14.\par
\textbf{Descripción:} El sistema deberá almacenar la información correspondiente a los procesos de votación y a su configuración.\par
\textbf{Datos específicos:} Identificador del proceso: entero positivo Organización: identificador de organización (RI-1) Título: 150 caracteres Descripción: 1000 caracteres Fecha de apertura: dd/mm/aaaa hh:mm Fecha de cierre: dd/mm/aaaa hh:mm (posterior a la de apertura) Estado: Borrador, Programado, Abierto, Cerrado, Escrutado o Archivado Re-votación permitida: Sí o No\par
\textbf{Comentarios:} El estado del proceso determina qué operaciones están permitidas (RF-08).\par

\subsection{RI-5: Información sobre papeletas}

\textbf{Objetivos asociados:} OBJ-2\par
\textbf{Requisitos asociados:} RF-06, RF-09, RF-10, RF-11\par
\textbf{Descripción:} El sistema deberá almacenar la información correspondiente a la papeleta de cada proceso, incluidas sus traducciones.\par
\textbf{Datos específicos:} Tipo de voto: Opción única o Sí/No/Abstención Pregunta: 300 caracteres Opciones (lista): texto de la opción (200 caracteres) y orden (entero positivo) Voto en blanco habilitado: Sí o No\par
\textbf{Comentarios:} Si falta una traducción, se muestra el texto en el idioma por defecto del proceso.\par

\subsection{RI-6: Información sobre votos (urna)}

\textbf{Objetivos asociados:} OBJ-3, OBJ-4\par
\textbf{Requisitos asociados:} RF-09, RF-10, RF-11, RF-13\par
\textbf{Descripción:} El sistema deberá almacenar los votos emitidos en cada proceso, sin ninguna relación con la identidad del votante.\par
\textbf{Datos específicos:} Identificador del voto: 36 caracteres (identificador aleatorio) Proceso: identificador del proceso (RI-4) Contenido del voto: opción elegida Fecha y hora de recepción: dd/mm/aaaa hh:mm:ss Canal: En línea Estado: Contabilizable o Sustituido\par
\textbf{Comentarios:} El voto no contiene ningún identificador ni dato de contacto del votante. Cuando un votante vuelve a votar (RF-11), el voto anterior pasa a Sustituido, sin borrarse. Solo se escrutan los votos Contabilizables. En el MVP el estado Sustituido no se genera (D-03): se conserva como referencia de diseño para RF-11 (P9), y activarlo requiere implementar la sustitución y validarla con pruebas de concurrencia y unicidad.\par

\subsection{RI-7: Información sobre participación}

\textbf{Objetivos asociados:} OBJ-3\par
\textbf{Requisitos asociados:} RF-08, RF-09, RF-11, RF-14\par
\textbf{Descripción:} El sistema deberá almacenar qué votantes han participado en cada proceso, sin indicar qué han votado.\par
\textbf{Datos específicos:} Proceso: identificador del proceso (RI-4) Votante: identificador de votante (RI-3) Ha votado: Sí o No Fecha y hora del último voto: dd/mm/aaaa hh:mm Canal del voto: En línea\par
\textbf{Comentarios:} Este registro se guarda separado de la urna (RI-6) y no enlaza con ningún voto concreto. Permite calcular la participación y evitar votos duplicados (RF-08/09).\par

\subsection{RI-8: Información sobre recibos de verificación}

\textbf{Objetivos asociados:} OBJ-3, OBJ-4\par
\textbf{Requisitos asociados:} RF-10, RF-11\par
\textbf{Descripción:} El sistema deberá almacenar la información correspondiente a los recibos que permiten al votante verificar su voto.\par
\textbf{Datos específicos:} Código del recibo: 64 caracteres hexadecimales Proceso: identificador del proceso (RI-4) Fecha y hora de emisión: dd/mm/aaaa hh:mm:ss Estado: Vigente o Sustituido\par
\textbf{Comentarios:} El recibo permite comprobar que el voto consta en la urna sin revelar su contenido. El recibo no se asocia a la identidad del votante.\par

\subsection{RI-9: Información sobre registro de auditoría de la urna}

\textbf{Objetivos asociados:} OBJ-4\par
\textbf{Requisitos asociados:} RF-07, RF-12\par
\textbf{Descripción:} El sistema deberá almacenar el registro encadenado de los eventos de la urna de cada proceso.\par
\textbf{Datos específicos:} Número de secuencia: entero positivo Proceso: identificador del proceso (RI-4) Tipo de evento: Apertura, Voto recibido, Voto sustituido, Cierre o Publicación de huella Fecha y hora: dd/mm/aaaa hh:mm:ss Huella del evento: 64 caracteres hexadecimales Huella del evento anterior: 64 caracteres hexadecimales Huella final publicada: 64 caracteres hexadecimales (solo tras el cierre)\par
\textbf{Comentarios:} Cada evento incluye la huella del anterior, de modo que cualquier alteración rompe la cadena y se detecta al verificarla (RF-12). Los eventos de voto no incluyen datos del votante.\par

\subsection{RI-10: Información sobre resultados y actas}

\textbf{Objetivos asociados:} OBJ-4\par
\textbf{Requisitos asociados:} RF-13, RF-15, RF-16\par
\textbf{Descripción:} El sistema deberá almacenar el resultado del escrutinio de cada proceso y el acta generada.\par
\textbf{Datos específicos:} Proceso: identificador del proceso (RI-4) Fecha del escrutinio: dd/mm/aaaa hh:mm Resultado por opción (lista): opción, votos (entero positivo) y porcentaje (número real con 2 cifras decimales) Votos en blanco: entero positivo Participación: número real (con 2 cifras decimales, en porcentaje) Huella del resultado: 64 caracteres hexadecimales Acta: fichero PDF Fecha de generación del acta: dd/mm/aaaa hh:mm\par
\textbf{Comentarios:} El resultado solo existe cuando el proceso está escrutado.\par

\subsection{RI-11: Información sobre notificaciones}

\textbf{Objetivos asociados:} OBJ-3\par
\textbf{Requisitos asociados:} RF-17\par
\textbf{Descripción:} El sistema deberá almacenar la información correspondiente a las notificaciones enviadas a los votantes y su estado de entrega.\par
\textbf{Datos específicos:} Identificador de notificación: entero positivo Proceso: identificador del proceso (RI-4) Votante: identificador de votante (RI-3) Tipo: Invitación, Recordatorio, Cierre Canal: Correo electrónico Plantilla: texto de hasta 1000 caracteres Fecha y hora de envío: dd/mm/aaaa hh:mm Estado: Pendiente, Entregada o Fallida Número de reintentos: entero de 1 cifra\par
\textbf{Comentarios:} Las plantillas y los canales los configura el Administrador de organización. Las notificaciones fallidas se pueden consultar (RF-17).\par

\subsection{RI-12: Información sobre acceso de votantes}

\textbf{Objetivos asociados:} OBJ-3\par
\textbf{Requisitos asociados:} RF-08, RF-17\par
\textbf{Descripción:} El sistema deberá almacenar la información necesaria para que cada votante acceda a la papeleta de un proceso abierto.\par
\textbf{Datos específicos:} Proceso: identificador del proceso (RI-4) Votante: identificador de votante (RI-3) Enlace personal: 64 caracteres (identificador aleatorio) Fecha de caducidad del enlace: dd/mm/aaaa hh:mm Código de verificación (OTP): 6 dígitos, se almacena solo su huella Fecha de caducidad del código: dd/mm/aaaa hh:mm Intentos fallidos: entero de 1 cifra Bloqueado hasta: dd/mm/aaaa hh:mm\par
\textbf{Comentarios:} El enlace personal identifica a un único votante en un único proceso. Tras varios intentos fallidos, el acceso se bloquea temporalmente.\par

\subsection{RI-13: Información sobre trazabilidad de acciones administrativas}

\textbf{Objetivos asociados:} OBJ-4\par
\textbf{Requisitos asociados:} RF-19\par
\textbf{Descripción:} El sistema deberá almacenar el registro de las acciones administrativas realizadas en la plataforma.\par
\textbf{Datos específicos:} Identificador de acción: entero positivo Fecha y hora: dd/mm/aaaa hh:mm:ss Usuario: identificador de usuario (RI-2) Organización: identificador de organización (RI-1) Tipo de acción: 50 caracteres Objeto afectado: 100 caracteres Detalle: 500 caracteres\par
\textbf{Comentarios:} El registro lo crea el Sistema automáticamente y no se puede modificar ni borrar. Es distinto del registro de auditoría de la urna (RI-10), que solo recoge eventos de la votación.\par

\section{Requisitos no funcionales}
\label{sec:rnf-detallados}

Los requisitos RNF-8 y RNF-10 se aplican de forma transversal a todos los
casos de uso (RF-01 a RF-19) y a todas las pantallas del sistema. No se
asocian a un caso de uso concreto.

\subsection{RNF-1: Accesibilidad WCAG 2.2 AA}

\textbf{Objetivos asociados:} OBJ-3\par
\textbf{Requisitos asociados:} RF-08, RF-09, RF-10, RF-11\par
\textbf{Descripción:} Las pantallas del votante deberán seguir el nivel AA de las pautas WCAG 2.2. Criterio de aceptación: cero incumplimientos críticos ni serios en las pantallas evaluadas, más revisión manual con teclado y lector de pantalla sin bloqueos en el flujo de acceso, voto y recibo.\par
\textbf{Comentarios:} Verificación: auditoría automática (por ejemplo, axe o Lighthouse) y revisión manual con teclado y lector de pantalla (H1a). La auditoría automática no certifica conformidad WCAG. Un incumplimiento serio puntual se documentará como incidencia y se decidirá si se corrige antes del cierre; el criterio no se relaja a priori. Criterio validado por el tutor. Sustituye al antiguo RF-18.\par

\subsection{RNF-2: Diseño adaptable y navegadores}

\textbf{Objetivos asociados:} OBJ-3\par
\textbf{Requisitos asociados:} RF-08, RF-09, RF-10, RF-11\par
\textbf{Descripción:} La interfaz deberá ser usable desde pantallas de 360 px de ancho, sin desplazamiento horizontal, y funcionar en las dos últimas versiones de Chrome, Firefox, Safari y Edge.\par
\textbf{Comentarios:} Verificación: pruebas en móvil, tableta y ordenador con cada navegador. Los valores (360 px, dos versiones) son propuestos. Sustituye al antiguo RF-18.\par

\subsection{RNF-3: Facilidad de uso del votante}

\textbf{Objetivos asociados:} OBJ-3\par
\textbf{Requisitos asociados:} RF-08, RF-09, RF-10, RF-11\par
\textbf{Descripción:} Las cinco personas no técnicas y sin formación previa que participen en la prueba deberán completar sin ayuda el acceso, el voto y la obtención del recibo en 3 minutos o menos.\par
\textbf{Comentarios:} Verificación: test de usabilidad con cinco personas no técnicas (H1b). Criterio del protocolo validado por el tutor. Con $n = 5$ no hay inferencia estadística: los resultados son evidencia descriptiva limitada al prototipo, la configuración y la población evaluados.\par

\subsection{RNF-4: Facilidad de uso del administrador}

\textbf{Objetivos asociados:} OBJ-1, OBJ-2\par
\textbf{Requisitos asociados:} RF-03, RF-06, RF-07\par
\textbf{Descripción:} Al menos 4 de 5 administradores de organización sin formación previa deberán crear, configurar y programar un proceso de opción única en 10 minutos o menos y sin ayuda. «Sin ayuda» significa sin intervención del moderador ni de otras personas; consultar la documentación de la herramienta sí está permitido.\par
\textbf{Comentarios:} Verificación: test de usabilidad con cinco personas con perfil de administrador. El valor de 10 minutos es propuesto. Con $n = 5$ no hay inferencia estadística: los resultados serán evidencia descriptiva limitada al prototipo, la configuración y la población evaluados. No hay hipótesis asociada en la sección~\ref{sec:plan-evaluacion}. La sesión del 03/11/2026 no cumplió las condiciones del criterio (población sin formación previa y tiempos medidos), por lo que RNF-4 se presenta como no demostrado en el capítulo de evaluación.\par

\subsection{RNF-5: Idiomas de la interfaz}

\textbf{Objetivos asociados:} OBJ-3\par
\textbf{Requisitos asociados:} RF-08, RF-09\par
\textbf{Descripción:} La interfaz deberá estar disponible en español e inglés, con todos los textos en ficheros de recursos, de modo que añadir un idioma no exija cambiar código. Cambiar de idioma no debe hacer perder los datos introducidos.\par
\textbf{Comentarios:} Verificación: revisión de código y prueba de cambio de idioma durante el flujo de voto.\par

\subsection{RNF-6: Claridad de la verificación del voto}

\textbf{Objetivos asociados:} OBJ-3, OBJ-4\par
\textbf{Requisitos asociados:} RF-10, RF-13\par
\textbf{Descripción:} La pantalla de verificación del recibo y del resultado deberá explicarse en lenguaje no técnico, sin términos criptográficos en el mensaje principal, y al menos 4 de las 5 personas no técnicas que participen en la prueba deberán interpretar correctamente si su voto consta en la urna.\par
\textbf{Comentarios:} Verificación: test con cinco personas no técnicas (H1b). Responde al objetivo de mejorar la usabilidad de la verificación para votantes no técnicos. Criterio del protocolo validado por el tutor. Con $n = 5$ no hay inferencia estadística: los resultados son evidencia descriptiva limitada al prototipo, la configuración y la población evaluados.\par

\subsection{RNF-7: Prevención de errores al votar}

\textbf{Objetivos asociados:} OBJ-3\par
\textbf{Requisitos asociados:} RF-09, RF-11.\par
\textbf{Descripción:} Antes de registrar un voto, el sistema deberá mostrar siempre un resumen de la selección y pedir confirmación explícita, y avisar de cualquier selección no válida antes de avanzar de pantalla.\par
\textbf{Comentarios:} Verificación: prueba funcional con todos los tipos de papeleta (RF-09), incluidas selecciones no válidas.\par

\subsection{RNF-8: Mensajes de error y realimentación}

\textbf{Objetivos asociados:} OBJ-1, OBJ-2, OBJ-3, OBJ-4\par
\textbf{Requisitos asociados:} Todos los requisitos funcionales (RF-01 a RF-19)\par
\textbf{Descripción:} Todo error mostrado al usuario deberá indicar con lenguaje claro qué ha ocurrido y cómo corregirlo, sin códigos técnicos. Toda acción que tarde más de 1 segundo deberá mostrar un indicador de progreso y toda acción completada deberá mostrar una confirmación visible.\par
\textbf{Comentarios:} Verificación: revisión de los mensajes de las excepciones de los casos de uso y prueba cronometrada de las acciones lentas. El valor de 1 segundo es propuesto.\par

\subsection{RNF-9: Orientación del votante durante el voto}

\textbf{Objetivos asociados:} OBJ-3\par
\textbf{Requisitos asociados:} RF-08, RF-09, RF-11.\par
\textbf{Descripción:} En cada paso del voto (acceso, selección, confirmación y recibo), el votante deberá ver en qué paso está y cuántos quedan, y en todo momento el estado del proceso y su hora de cierre.\par
\textbf{Comentarios:} Verificación: revisión de las pantallas del flujo de voto y test de usabilidad (RNF-3).\par

\subsection{RNF-10: Consistencia de la interfaz}

\textbf{Objetivos asociados:} OBJ-1, OBJ-2, OBJ-3, OBJ-4\par
\textbf{Requisitos asociados:} RF-01 a RF-19\par
\textbf{Descripción:} La interfaz deberá usar la misma terminología, la misma ubicación de los elementos de navegación y los mismos componentes en todas las pantallas, sin problemas graves de consistencia en una revisión heurística.\par
\textbf{Comentarios:} Verificación: revisión heurística (heurísticas de Nielsen) por al menos 2 evaluadores, sin problemas de gravedad alta. Los valores son propuestos.\par

\subsection{RNF-11: Secreto del voto}

\textbf{Objetivos asociados:} OBJ-3, OBJ-4\par
\textbf{Requisitos asociados:} RF-09, RF-10, RF-11, RF-13\par
\textbf{Descripción:} Ninguna tabla enlaza voto e identidad. La urna (RI-6) y la participación (RI-7) están separadas.\par
\textbf{Comentarios:} Verificación: revisión del modelo de datos y prueba de no correlación.\par

\subsection{RNF-12: Unicidad}

\textbf{Objetivos asociados:} OBJ-3\par
\textbf{Requisitos asociados:} RF-09, RF-11\par
\textbf{Descripción:} El sistema deberá garantizar que solo un voto sea contabilizado por votante, incluso cuando se produzcan envíos simultáneos.\par
\textbf{Comentarios:} Verificación: prueba de envíos concurrentes del mismo votante para comprobar que únicamente un voto queda en estado Contabilizable.\par

\subsection{RNF-13: Integridad y trazabilidad}

\textbf{Objetivos asociados:} OBJ-4\par
\textbf{Requisitos asociados:} RF-12, RF-19\par
\textbf{Descripción:} El registro usa hash encadenado y cualquier alteración se detecta. Toda acción administrativa queda registrada y no es editable.\par
\textbf{Comentarios:} Verificación: prueba de alteración del registro y auditoría.\par

\subsection{RNF-14: Seguridad de acceso}

\textbf{Objetivos asociados:} OBJ-1, OBJ-3\par
\textbf{Requisitos asociados:} RF-04, RF-08\par
\textbf{Descripción:} Contraseñas con hash y sal. Bloqueo tras 5 intentos fallidos. Segundo factor. HTTPS. Sesión con caducidad.\par
\textbf{Comentarios:} Verificación: pruebas de seguridad. Valores propuestos.\par

\subsection{RNF-15: Aislamiento entre organizaciones}

\textbf{Objetivos asociados:} OBJ-1\par
\textbf{Requisitos asociados:} RF-01\par
\textbf{Descripción:} Ningún usuario accede a datos de otra organización, comprobado con pruebas.\par
\textbf{Comentarios:} Verificación: pruebas de acceso cruzado entre tenants.\par

\subsection{RNF-16: Protección de datos (RGPD)}

\textbf{Objetivos asociados:} OBJ-1\par
\textbf{Requisitos asociados:} RF-18\par
\textbf{Descripción:} El sistema deberá limitar la recogida de datos personales a los estrictamente necesarios para la prestación del servicio y deberá permitir su eliminación cuando corresponda conforme a la normativa aplicable.\par
\textbf{Comentarios:} Verificación: revisión del modelo de datos y prueba del flujo de eliminación de datos personales.\par

\subsection{RNF-17: Rendimiento}

\textbf{Objetivos asociados:} OBJ-3\par
\textbf{Requisitos asociados:} RF-09, RF-13\par
\textbf{Descripción:} El sistema deberá soportar al menos 500 votantes simultáneos durante la emisión de votos y deberá responder en menos de 2 segundos en al menos el 95 \% de las peticiones realizadas durante dicha carga.\par
\textbf{Comentarios:} Verificación: prueba de carga con al menos 500 usuarios simultáneos y medición del tiempo de respuesta. Los valores establecidos son propuestos y deberán validarse con el cliente.\par

\subsection{RNF-18: Portabilidad (PWA instalable)}

\textbf{Objetivos asociados:} OBJ-3\par
\textbf{Requisitos asociados:} RF-08, RF-09\par
\textbf{Descripción:} La aplicación deberá ser instalable como PWA en dispositivos móviles.\par
\textbf{Comentarios:} Prioridad C. Verificación: instalación en Android e iOS.\par

